\documentclass[11pt,a4paper]{article}
\usepackage[margin=25mm]{geometry}
\usepackage{iftex}
\ifPDFTeX\usepackage[T1]{fontenc}\usepackage[utf8]{inputenc}\fi
\usepackage{amsmath,amssymb,graphicx,booktabs,array}
\usepackage[authoryear,round]{natbib}
\usepackage{caption,placeins}
\usepackage[unicode,breaklinks=true,hidelinks]{hyperref}
\setlength{\parindent}{0pt}\setlength{\parskip}{5pt}
\setlength{\emergencystretch}{3em}\setlength{\tabcolsep}{3pt}
\clubpenalty=10000\widowpenalty=10000
\captionsetup{font=small,labelfont=bf}
\newcounter{proposition}
\begin{document}
\begin{center}\Large\bfseries Hard to Beat: Historical Rules versus Conditional Quantile Models for Peak Net-Position Limits\end{center}

\begin{abstract}Risk limits on dealer positions concern the largest magnitude reached before a review horizon ends, rather than the position at its endpoint. We compare historical and conditional quantile forecasts of sampled absolute-position peaks for an FX market maker and weekly primary-dealer position aggregates. Simple historical rules prove difficult to beat. At the shorter horizon, long-history FX regression lowers matched loss by 14.8\% and 16.6\% in the previously reported later and earlier windows. Its later advantage survives a delayed adaptive coverage adjustment (ACI). Before correction, however, regression breaches more frequently, while its ranking against a shorter-history rolling rule changes across periods. After stronger regularization, regression still loses 34 of 48 additional-panel comparisons. We formalize the known-floor target and apply a classical endpoint-to-peak bound under explicit assumptions. The results support joint validation of accuracy, breach behavior, and stability across reference rules, periods, and related series. A controlled simulation rejects a universal narrow-regime explanation, while aggregation and the overlapping backcast further qualify the empirical pattern.\end{abstract}

\section{Introduction}\label{sec:intro}

A dealer's recorded net position can finish inside its limit after exceeding it earlier in the monitoring period, as Figure~\ref{fig:example} illustrates. The forecasting question concerns the largest sampled exposure before the next review, rather than the endpoint alone. This changes the response and when its training label becomes available. We ask whether additional conditioning improves that peak forecast once historical information is matched, and whether the gain survives alternative references and breach assessment.


Simple historical rules provide a demanding reference because they already incorporate information about the current position and its recent path behavior. A conditional model must improve on that information while estimating additional relationships from a limited and changing sample of past peaks. More predictors can help distinguish current conditions, but they also increase sensitivity to the historical sample and the strength of regularization. We examine these competing effects through matched histories, alternative reference rules, and repeated chronological comparisons.


\begin{figure}[htbp]\centering\includegraphics[width=.94\linewidth]{figures/v7-endpoint-peak-example.pdf}
\caption{A selected Treasury-bill path whose endpoint returns within the historical report after an intervening breach. The rule chooses the earliest eligible original evaluation window. Endpoint containment can conceal a peak breach.}\label{fig:example}
\end{figure}

We study an hourly FX market-maker record and weekly U.S. primary-dealer position aggregates, asking where conditioning improves peak forecasts. In the original ordered comparisons, historical rules, regression and local quantiles compete for the same peak using issue-known observations. These comparisons reveal whether gains persist across reference rules and periods, and whether better scores also mean fewer breaches.


The clearest repeated gains arise for long-history FX regression at the shorter forecast horizon, rather than across every conditional specification. Against a historical headroom rule with the same history, it lowers later-period loss by 14.8\%. The earlier-period reduction is 16.6\%, and the delayed coverage adjustment preserves the later advantage. Changing the comparison to the shorter-history rolling reference reverses the direction of the earlier-period result, while weekly comparisons frequently favor historical rules.


A lower quantile score also leaves a separate operational question about how often the predicted limit is crossed during the forecast period. The later matched FX comparison records 376 regression breaches, compared with 360 for its historical reference. For the separately named comparisons, Murphy diagrams plot mean elementary quantile scores across thresholds, revealing where consistent scoring rules give different rankings. Error-weight comparisons show how preferences between fixed forecasts depend on the relative penalties on underprediction and overprediction. For managers setting monitoring limits, the practical implication is to evaluate prediction errors and exceedance frequency together before selecting a more complex forecast.


\textbf{C1.} We formalize the peak-position quantile as a path-dependent forecast with a known current-position floor. We record an elementary loss-improving floor projection and the true-quantile strict-tail bound, then apply the classical Lévy inequality under independent symmetric increments. Separately fitted short-history FX regression forecasts cross across horizons on 640 evaluation origins. Wisniewski et al. (2020) forecast signed endpoint intervals rather than a sampled absolute inventory peak. The portfolio-loss maximum of Boudoukh et al. (2004) and payment-deficit peak of van den Hoogen (2017) provide distinct path-risk precedents. Section~\ref{sec:target} provides the definitions, propositions and crossing diagnostic, with proofs and simulation qualifications in the Supplement.

\textbf{C2.} We find that historical rules win most weekly comparisons, while the original long-history FX gain depends on the chosen reference. We compare the same peak target for one FX record and thirteen weekly position-category aggregates with matched histories and update schedules. Even after extending the regression penalty menu, historical rules win in 34 of 48 additional-panel comparisons. Where Duffie et al. (2023) examine capacity, illiquidity and PnL quantiles, and Wisniewski et al. (2020) forecast one market maker's endpoints, we compare these peak-position tasks. Section~\ref{sec:results} reports the gains and failures, and Table~\ref{tab:map} organizes the observed settings and validation questions. The controlled study rejects the proposed regime restriction, while aggregation supplies a same-dealer contrast without isolating frequency.

\textbf{C3.} We show how peak-limit rankings change with the historical reference and breach assessment in these observed forecasting tasks. We compare same-history and shorter-history rolling controls across three windows with previously examined outcomes, using established score diagrams and error-weight preferences. For the distinct long-horizon legacy comparison, the preferred report reverses at a penalty ratio of about 20.6:1. Where the endpoint lineage evaluates intervals and Ehm et al. (2016) develop general score comparisons, we connect these diagnostics to monitoring limits with completed path labels. Sections~\ref{sec:evaluation}--\ref{sec:implications} support a four-part checklist: use matured labels, separate history and model changes, report scores with breaches, and check stability across periods and related series.


The remaining sections define the peak target, describe the data and forecasts, and compare accuracy, breaches and stability before drawing implications for model validation. Complexity earns its place through a stable comparative advantage on the quantity that the limit actually monitors.


\FloatBarrier

\section{Related literature}\label{sec:literature}

Net-position forecasting already has a close empirical lineage in the market-making literature, where Wisniewski and colleagues predict signed next-hour endpoint intervals. Their methods combine historical positions, calendar information, neighboring observations, and normalized prediction intervals, with sequential training, calibration, and evaluation \citep{wisniewski2020,wisniewski2024}. They also report scores and coverage tests together, providing a direct precedent for the joint evaluation of interval forecasts. Our response instead requires the complete future path segment and includes the current absolute position as a known lower bound. This distinction creates different training labels and different limits, even when the numerical forecasting engine follows the same broad family.


Financial path-risk studies establish that a within-period extreme can matter more than the value observed at the end of the period. MaxVaR evaluates the worst cumulative portfolio return under a fixed portfolio and diffusion assumptions, including a conversion from endpoint probabilities to path risk \citep{boudoukh2004}. The payment-liquidity thesis studies conditional maximum deficits using a different path-risk target and confidential transactions \citep{hoogen2017}. Related settlement studies examine tail liquidity needs and collateral requirements in their respective payment systems \citep{petrunia2019,arjani2020}. Duffie and colleagues analyze dealer capacity and Treasury-market functioning, with additional out-of-sample PnL quantile comparisons in their appendix \citep{duffie2023}. Our empirical contribution concerns matched forecasts of the present peak-position target across the FX record and related public position aggregates.


Quantile estimation and evaluation supply established tools for comparing conditional peak forecasts against historical rules on a shared response. Regression quantiles directly target asymmetric forecast errors through the associated loss function, providing the basic statistical building block \citep{koenker1978}. Conditional autoregressive quantiles and quantile forests provide additional conditioning structures for different forecasting applications \citep{engle2004,meinshausen2006}. Consistent scoring rules and elementary-score representations explain why rankings can vary across distinct evaluation questions \citep{gneiting2011,ehm2016}. We connect these tools to mature path labels, matched historical controls, and the stability of observed score and breach differences. Table~\ref{tab:literature} shows how the peak response, mature training labels and matched reference rules distinguish the present forecasting comparison.


\begin{table}[htbp]\centering\footnotesize\uchyph=1\hyphenation{con-fi-den-tial con-di-tion-al re-gres-sion com-plet-ed el-e-men-ta-ry ex-pand-ing}
\caption{Closest literature and the present forecasting task. n.r. marks a dimension unreported in the consulted task or passage. The distinction is a common peak-position target and its validation design.}\label{tab:literature}
\begin{tabular}{>{\raggedright\arraybackslash}p{2.4cm}>{\raggedright\arraybackslash}p{2.1cm}>{\raggedright\arraybackslash}p{1.5cm}>{\raggedright\arraybackslash}p{1.6cm}>{\raggedright\arraybackslash}p{1.6cm}>{\raggedright\arraybackslash}p{1.7cm}>{\raggedright\arraybackslash}p{1.7cm}}\toprule
Study & Response & Data unit & Peak floor & Label timing & Evaluation & Source scope\\\midrule
\shortstack[l]{Wisniewski\\et al. (2020)} & Signed endpoint interval & One FX market maker & n.r. & Ordered train/\allowbreak calibrate/\allowbreak test & Scores and hit tests & One source family\\
\shortstack[l]{Boudoukh\\et al. (2004)} & Worst portfolio return & Stylized fixed portfolio & Different initial-value anchor & Analytic diffusion & Tail ratios and monitoring & Parameter cases\\
\shortstack[l]{van den\\Hoogen (2017)} & Cumulative payment-deficit peak & Con\-fi\-den\-tial bank payments & Absolute-position floor n.r. & Path simulation & EVT methodology; results partly confidential & Bank payment task\\
\shortstack[l]{Duffie\\et al. (2023)} & Capacity, illiquidity, PnL quantiles & Dealer/\allowbreak market sources & n.r. & Expanding PnL fits & Capacity and PnL forecasts & Several sources\\
\shortstack[l]{Engle and\\Manganelli\\(2004)} & Return quantile & Asset returns & n.r. & Lagged predictors & Quantile scores and DQ & Several returns\\
\shortstack[l]{Ehm et al.\\(2016)} & Quantile/\allowbreak expectile comparison & Multiple forecast examples & n.r. & Original example designs & El\-e\-men\-tary-score curves & Multiple tasks\\
\shortstack[l]{Meinshausen\\(2006)} & Conditional response quantile & Regression examples & n.r. & Original regression splits & Conditional distributions & Several datasets\\
This study & Sampled absolute-position peak & One FX record; weekly totals & Observed current magnitude & Completed forward labels & Scores, breaches, coherence & Related aggregate tasks\\
\bottomrule\end{tabular}
\par\smallskip\begin{minipage}{.98\linewidth}\footnotesize The institutional payment-thesis reading is partial; n.r. is not evidence of absence. The earlier FX work and this paper share a source family, without verified input-byte identity.\end{minipage}
\end{table}

\FloatBarrier

\section{The forecast target}\label{sec:target}

The forecast target is the largest sampled absolute position between the issue time and the end of the stated forecast horizon. Let \(N_t\) denote the reported signed position and let \(\mathcal F_t\) contain the information available at the forecast origin. Define the peak \(Y_{t,H}\), its observed floor \(B_t\), and the corresponding nonnegative headroom \(Z_{t,H}\) by


\[
Y_{t,H}=\max_{0\le h\le H}|N_{t+h}|,\qquad B_t=|N_t|,\qquad Z_{t,H}=Y_{t,H}-B_t\ge0.
\]

The inclusion of the issue position matters because a feasible peak limit must already cover the exposure observed at the forecast origin. Figure~\ref{fig:example} shows a real window selected by a fixed chronological rule, with an endpoint inside the limit and an earlier peak outside it. A successful endpoint interval can therefore answer a different risk question from a successful report for the full sampled peak. We evaluate the latter through pinball loss and count a breach precisely when the observed peak strictly exceeds the reported limit. This report monitors sampled net inventory; its known floor neither measures settlement cash needs nor establishes an optimal operating limit.


\[
\rho_\tau(u)=\tau u_+ +(1-\tau)(-u)_+,\qquad I_{t,H}=\mathbf 1\{Y_{t,H}>C_{t,H}\}.
\]

\refstepcounter{proposition}\noindent\textbf{Proposition~\theproposition (known-floor projection).}\label{prop:floor} Suppose \(Y\ge B\) almost surely and the issue information contains \(B\). Projecting any forecast \(c\) to \(\max\{B,c\}\) weakly improves pinball loss for every possible outcome, with an exact gain below the floor. If \(q\) is the true conditional lower inverse quantile at level \(1-\alpha\), its strict exceedance probability is at most \(\alpha\). At an atom of the response distribution, that probability can be smaller than the nominal tail probability.


\[
\rho_\tau(Y-c)-\rho_\tau(Y-\max\{B,c\})=\tau(B-c)_+,\qquad \Pr(Y>q\mid\mathcal F_t)\le\alpha.
\]

\refstepcounter{proposition}\noindent\textbf{Proposition~\theproposition (classical endpoint-to-peak bound).}\label{prop:levy} Let future increments be independent and individually symmetric about zero, with partial sums \(S_h\). Fix a starting position \(x\), set \(B=|x|\), and let \(q\) be the population \((1-\alpha/2)\) quantile of \(|S_H|\). The triangle inequality and the classical scalar L\'evy maximal inequality then imply the following bound, including strict inequalities at atoms \citep{sussman2023}. Neither identical increment distributions nor finite increment moments are required for this population statement.


\[
\Pr\!\left(\max_{0\le h\le H}|x+S_h|>B+q\right)\le\Pr\!\left(\max_h|S_h|>q\right)\le2\Pr(|S_H|>q)\le\alpha.
\]

The Supplement gives both proofs and distinguishes the population endpoint quantile from the empirical endpoint estimate used by the corresponding benchmark. A conditional use of Proposition~\ref{prop:levy} requires the independence, symmetry, and endpoint-quantile assumptions under the law conditional on issue information. Its simulation comparisons separate mean-reverting position levels from serially correlated increments, since those constructions have different dependence in their future paths. This separation makes the mathematical premises explicit before the empirical benchmark is judged by its observed losses and breaches.


Forecasts for longer horizons should also respect the nesting of the underlying path maxima, which makes the true horizon quantiles ordered. Separate fits can violate that order even though their observed targets remain correctly nested at every evaluated forecast origin. The short-history FX regression reports cross across horizons on 640 origins, illustrating a coherence problem distinct from the score of each marginal forecast. The target definition therefore determines the floor, label maturity, breach event, and horizon relationships that a peak-limit comparison must preserve.


\FloatBarrier

\section{Data}\label{sec:data}

The FX application uses the hourly member of the publicly described Net Position Time Series Dataset, representing a market maker's net client-flow position \citep{dataset,wisniewski2020}. The accepted source contains 18576 records, with positions measured in the source's USD units and ordered by its observation clock. We use the same series for the later evaluation and the earlier chronological comparison, with their separate evaluation indices preserved throughout. Within each comparison, every candidate forecast is scored against the same sampled peak and the same current-position floor.


The weekly application uses Primary Dealer Statistics published by the Federal Reserve Bank of New York \citep{nyfeddata}. Its reporting definitions follow the government-securities dealer framework and the relevant position-accounting instructions \citep{fr2004_2013,fr2004_2024}. The original Treasury-bill aggregate has 706 observations, with values expressed in USD millions and indexed by their economic dates. Additional position categories use the same complete weekly date grid after a metadata and calendar eligibility check conducted before their model evaluation. These categories cover different reported position components and are fitted separately rather than pooled into a common monetary forecasting sample.


The FX record describes one market maker, while the weekly series combine signed positions across primary dealers. The economic-date convention identifies where the observations enter the downloaded series, while the reporting definitions identify the positions being aggregated. Each analysis uses its own units, training history, forecast horizon, and issue grid, with reference rules matched inside the relevant comparison. The data design makes frequency and measurement differences visible instead of treating every sequence as an interchangeable dealer history.

Daily aggregation of the same dealer produces 930 nonempty source-naive calendar days and 501 common evaluation issues. Each day contributes its last record, including partial boundary days, and the report clock follows the end of that day. The timezone and historical release vintage remain unknown; this last-record proxy is distinct from a verified market close or intraday maximum.


\FloatBarrier

\section{Forecasting methods}\label{sec:methods}

The primary forecasts target the conditional quantile at \(\tau=19/20\), with its upper tail defining the strict-breach event. Its nineteen-to-one error preference is hypothetical and is separate from measured funding, trading or settlement costs. The two horizons span \(H\in\{6,24\}\) subsequent source records, and each peak also includes the position observed at issue time. This common target keeps the comparisons aligned while allowing the historical pool and conditioning structure to differ between the named procedures.


Historical rules serve as substantive competitors because they can already adapt their limits to the position observed at the forecast origin. The historical headroom rule adds an empirical quantile of completed past headroom to the current absolute position, while the translated-path rule reanchors historical signed paths. The historical-peak rule uses completed absolute peaks directly, and the endpoint-derived rule adds an empirical absolute-displacement quantile to the observed floor. Table~\ref{tab:methods} gives their definitions, update schedules, and the parameter menus used by the conditional methods.


Conditional quantile regression links completed peak headroom to issue-known position histories and calendar information, using training-based transformations and ridge regularization. The original predictor set contains 33 features before its intercept, with all transformations fitted inside the chronological training sample. The parsimonious variant concentrates on a smaller state description, while the expanded ridge comparison changes the penalty menu with the richer predictor set held fixed. These comparisons distinguish changes in the conditioning information from changes in the strength of regularization rather than attributing every gain to feature reduction.


Local quantile methods condition on neighboring historical states, with a scale-adjusted variant using training information to normalize the comparison. The filtered-path benchmark recursively simulates signed positions using GARCH-standardized empirical innovations drawn independently with replacement \citep{hull1998,barone1999}. The other benchmarks estimate a quantile from response distributions in tree leaves or fit tails of volatility-scaled peak headroom \citep{meinshausen2006,mcneilfrey2000}. These are the existing fitted specifications, with every complete result and numerical or structural failure documented in the Supplement. Their different conditioning assumptions make the actual executed procedure, rather than its family name alone, the relevant object for comparison.


The further robustness comparisons address calibration, forecast combination, and conservative selection while keeping the peak response and chronological information restrictions explicit. A headroom version of the symmetric absolute-value quantile recursion follows the CAViaR family, and finite delayed calibration wrappers update only after outcomes mature \citep{engle2004,gibbs2021}. Fixed mean and median combinations aggregate the specified forecasts, while a paired-difference one-standard-error rule favors simpler reports within its stated validation tolerance. Every method earns its comparison against the historical headroom rule through its actual score and breach outcomes on the shared evaluation origins.


\begin{table}[htbp]\centering\small
\caption{Forecast definitions and executed menus. Codes occur here to connect with the result files. Matched histories and update schedules make the reference explicit.}\label{tab:methods}
\begin{tabular}{>{\raggedright\arraybackslash}p{2.9cm}>{\raggedright\arraybackslash}p{6.0cm}>{\raggedright\arraybackslash}p{5.3cm}}\toprule
Method / code & Forecast definition & Executed menu / setting\\\midrule
Historical headroom (F) & Current floor plus a quantile of completed peak headrooms & FX histories \{720,2880\}; weekly \{52,260\}\\
Translated path (S) & Translate completed signed paths to the current position before taking their peaks & Same completed-label pools; fixed-window and rolling-window controls\\
Historical peak (K) & Historical absolute-peak quantile projected to the current floor & Same histories and update controls\\
Quantile regression (QR) & Training-scaled headroom regression with ridge slopes and an unpenalized intercept & Original penalties $\{0.001,0.01,0.1\}$; expanded menu in Supplement\\
Local quantiles (KNN) & Neighbor headroom quantile, with a separate increment-scale variant & Training-only transforms and fixed neighbor menus; Supplement\\
Parsimonious regression (QR4) & Current signed/magnitude positions and short/long increment RMS & Separate five-point penalty menu; common-menu comparison in Supplement\\
Endpoint plug-in (L) & Floor plus a high absolute-displacement endpoint quantile & Proposition~\ref{prop:levy}; empirical plug-in evaluated separately\\
Further benchmarks & Autoregressive replay; filtered paths; peak-tail fit; quantile forest & Existing bounded specifications, complete failures included in Supplement\\
Coverage and combinations & Delayed projected ACI-style wrappers; headroom SAV; fixed forecast combinations & Original-source wrapper/SAV inventory and broader combination inventory are separate\\
\bottomrule\end{tabular}
\par\smallskip\begin{minipage}{.98\linewidth}\footnotesize Training histories count completed labels. Fixed-window models refit every 24 FX records or 13 weekly records. Quantile and normalization conventions are detailed in the Supplement; this table does not expand any executed menu.\end{minipage}
\end{table}

\FloatBarrier

\section{Evaluation design}\label{sec:evaluation}

The forecasting clock admits a training label only after the entire response horizon has ended, preventing an unfinished peak from entering estimation. At refit origin \(r\), a historical origin \(i\) is eligible only when \(i+H\le r\), and every selected window uses the same completed-label convention. Predictors use issue-known observations, and training-only transformations remain fixed until the next scheduled refit rather than learning from future evaluation outcomes. The original chronological comparisons select parameters in the preceding validation segment and evaluate the complete specified issue grid. The expanded backcast instead retains penalties selected later, although each coefficient fit and transformation still uses only mature training labels.


Matched comparisons hold the completed-response history and refit schedule fixed when evaluating the regression against the historical headroom rule. The original FX primary instead uses the shorter-history rolling historical reference, making it a distinct comparison from the same-history benchmark. The Treasury-bill primary compares the longer-history regression with the historical headroom rule on a shared history and update schedule. Tables~\ref{tab:fx} and~\ref{tab:weekly} preserve these original comparisons alongside the matched, additional-panel, and earlier-period results.


We evaluate forecast accuracy with mean pinball loss and report strict breaches separately, including their frequency and the associated overrun amounts. Paired loss differences use heteroskedasticity and autocorrelation consistent covariance estimates, with the horizon represented in the bandwidth choice \citep{dm1995,nw1987}. The existing calculations also report the small-sample adjustment and the stated multiple-comparison families, with each family's role identified explicitly \citep{hln1997,holm1979}. Coverage diagnostics use a fixed spaced grid and distinguish unconditional coverage from a transition likelihood conditioned on its initial hit \citep{kupiec1995,christ1998}.


\subsection{Research design and transparency}

The original source-specific primaries retain their original roles, outcomes, and reference procedures throughout this paper. The conditional extension followed inspection of the original FX evaluation archive; later calibration, combination, state, quantile-level, and selection analyses are post-result and exploratory. Their designs and execution records preserve the chronology of each addition, including failed candidates and full comparison families rather than selected favorable subsets. The expanded prefix reuses familiar source observations, overlaps the earlier backcast, and follows previous rolling-validation scoring of its target paths. The new extensions were fixed before their own calculations within an already inspected research program. Their tests and the writing-stage HAC comparisons remain descriptive rather than independent or prospective confirmation.


The Supplement provides the complete prior and later comparison inventories, the exact evaluation clocks, numerical checks, and all incomplete or unsuccessful procedures. We report sensitivity to bandwidth, adjustment family, target level, and state definition as separate diagnostic questions within their specified scopes. The subsequent interpretation follows the same sequence: establish the reference and sample, compare losses, examine breaches, and assess stability across periods. A credible peak-limit gain must survive this sequence with its benchmark and its information set stated clearly.


\FloatBarrier

\section{Results}\label{sec:results}

\subsection{Where conditional models improve}

Long-history FX regression provides the clearest repeated improvement against a historical headroom rule using the same completed history and update schedule. It lowers matched loss by 14.8\% in the later period and 16.6\% in the earlier period. The same-history comparison holds the amount of training information fixed, whereas the shorter-history rolling reference also changes the available history. The matched gains in Table~\ref{tab:fx} therefore isolate a different comparison from the original rolling-reference result.


The delayed ACI adjustment preserves the matched FX score advantage while bringing the two observed breach frequencies close together. Adjusted regression lowers mean loss by 15.0\% against the adjusted historical reference on their common evaluation grid. Their strict breach counts are 329 and 328, respectively. This comparison shows that the score gain survives the adjustment, with nearly equal breach frequencies for these two implemented procedures.


The long-history FX gain also persists across the evaluated position-increment states and at the lower target quantile. The unchanged primary has a negative mean loss difference in every training-defined RMS group, with the group results documented in the Supplement. At the lower target, regression lowers loss by 10.1\% against matched headroom and by 9.2\% against the original rolling reference. These results broaden the description of that configuration while leaving the separate upper-target validation failures visible in the full record.


\begin{table}[htbp]\centering\small
\caption{Long-history short-horizon FX regression under two references and periods. Negative percentages favor regression. Matched gains recur while the rolling-reference direction changes.}\label{tab:fx}
\begin{tabular}{>{\raggedright\arraybackslash}p{2.6cm}>{\raggedright\arraybackslash}p{4.4cm}>{\raggedright\arraybackslash}p{2.5cm}>{\raggedright\arraybackslash}p{1.5cm}>{\raggedright\arraybackslash}p{3.2cm}}\toprule
Period / procedure & Reference & Loss change (\%) & Issues & Strict breaches (regression / reference)\\\midrule
Later / original primary & Shorter-history rolling headroom & -12.66 & 6625 & 376 / 332\\
Later / matched & Same-history fixed-window headroom & -14.77 & 6625 & 376 / 360\\
Earlier / original structure & Shorter-history rolling headroom & +3.17 & 1440 & See complete earlier-period tables\\
Earlier / matched & Same-history fixed-window headroom & -16.62 & 1440 & 20 / 16\\
Later / coverage-wrapped & Same-history wrapped headroom & -15.00 & 6625 & 329 / 328\\
\bottomrule\end{tabular}
\par\smallskip\begin{minipage}{.98\linewidth}\footnotesize Regression uses history 2880 and horizon 6; the original rolling reference uses 720. The earlier matched comparison is the original three-penalty regression, not the later one-standard-error selection rule. The wrapper comparison is a separate finite projected delayed ACI-style exercise.\end{minipage}
\end{table}

\subsection{Where conditional models lose}

The advantage changes when the reference changes, even for the same long-history FX regression report at the shorter horizon. Against the original shorter-history rolling rule, it lowers later-period loss by 12.7\% but raises earlier-period loss by 3.2\%. The earlier matched improvement therefore coexists with an adverse original-primary comparison, rather than reproducing the original gain against every historical rule. The reversal in Table~\ref{tab:fx} shows that a matched-history gain can coexist with failure against the simpler operational reference.


Short histories and longer horizons produce additional adverse outcomes that prevent a broad ranking of regression above historical forecasting rules. The short-history FX regression at the longer horizon raises mean loss by 31.5\% relative to its matched historical headroom rule. Its breach counts are 924 for regression and 427 for the reference. Other long-horizon comparisons show smaller score gains alongside more breaches, so a positive result at the shorter horizon cannot stand in for the full map.


The original Treasury-bill primary loses on both score and breach count. Its loss increase is 50.8\%, with 54 regression breaches on the evaluation grid. The historical reference records 28 breaches, as Table~\ref{tab:weekly} shows together with the corresponding fixed-reference inference. The earlier shorter-history structural analogue also loses, so the weekly evidence challenges transport of the successful FX configuration to these aggregate tasks.


Heavier regularization improves the additional weekly panel substantially but leaves historical rules ahead in most of its evaluated cells. The richer regression's median excess loss falls from 92.6\% to 9.2\% after expansion of its ridge menu. Even then, it loses 34 of 48 comparisons against matched historical headroom rules. The earlier-period one-standard-error exercise misses its success criteria, reinforcing the need to examine stability rather than selecting a favorable regularization summary.


Simpler conditional descriptions and forecast combinations also produce mixed outcomes rather than a uniform solution to the weekly regression losses. The fixed median combination lowers loss against the bill historical reference by 4.0\%. The fixed mean combination instead raises loss by 2.8\%, and the headroom quantile recursion also performs worse than that reference. Table~\ref{tab:weekly} and the Supplement preserve these comparisons together with the additional model failures, including every unfavorable completed result.


\begin{table}[htbp]\centering\small
\caption{Weekly evidence across the original primary, added panel and earlier-period comparisons. Counts refer to related tasks rather than independent dealers. Stronger regularization narrows losses while leaving substantial failures.}\label{tab:weekly}
\begin{tabular}{>{\raggedright\arraybackslash}p{4.3cm}>{\raggedright\arraybackslash}p{4.7cm}>{\raggedright\arraybackslash}p{2.7cm}>{\raggedright\arraybackslash}p{2.5cm}}\toprule
Comparison set & Procedure & Loss change / median (\%) & Higher-loss cells\\\midrule
Original Treasury-bill primary & Original regression, long history & +50.79 & Original primary\\
Earlier Treasury-bill analogue & Short-history structural analogue & +144.30 & Different history\\
Additional panel & Original full regression & +92.63 & 45 / 48\\
Additional panel & Four-predictor regression & +9.91 & 37 / 48\\
Additional panel & Full regression, expanded penalties & +9.22 & 34 / 48\\
Additional panel & Ordinary empirical headroom quantile & +0.50 & 30 / 48\\
Earlier weekly groups & Validation-minimum expanded rule & +4.15 & 26 groups\\
Earlier weekly groups & One-standard-error rule & +5.37 & 26 groups\\
\bottomrule\end{tabular}
\par\smallskip\begin{minipage}{.98\linewidth}\footnotesize The original bill primary has 54 versus 28 strict breaches among 277 issues. Earlier short-history bill results are a structural analogue, not an exact replay of that primary. The one-standard-error worst excess is +181.39\%; all its prewritten success criteria remain unmet. Full panel and selection tables are linked from the Supplement and EXTENDED\_RESULTS.md.\end{minipage}
\end{table}

\subsection{Scores and breach behavior answer different questions}

The matched FX gain comes with more breaches before calibration, showing why pinball accuracy and exceedance frequency deserve separate reporting. Regression breaches 376 times, compared with 360 for the matched historical headroom rule. Pinball loss measures the sizes of underprediction and overprediction; a strict breach records whether the realized peak exceeds its forecast. Managers comparing monitoring limits therefore need both measures, rather than interpreting a smaller average score as a direct reduction in breach occasions.


Murphy diagrams reveal that forecast rankings can depend on the scoring rule even when the quantile level and reports stay fixed. Figure~\ref{fig:murphyfx} displays crossing curves for the original FX primary, and Figure~\ref{fig:murphybill} presents the distinct Treasury comparison. A separate error-weight analysis finds a preference reversal near 20.6:1 for the long-horizon legacy FX report against the rolling translated-path reference. The short-horizon FX primary and matched comparisons have lower mean overrun and unused allowance together, so that particular reversal belongs to the long-horizon comparison.


Table~\ref{tab:diagnostics} places the working-reference loss tests alongside the fixed-grid coverage results, preserving their distinct null hypotheses and comparison families. A loss-difference test concerns average relative accuracy, while coverage tests concern the frequency and temporal pattern of forecast breaches. State diagnostics add another comparison between conditional predictive ability and differences across the specified RMS groups \citep{gw2006}. The combined evidence supports a narrow score advantage in the successful FX setting and a broader requirement to assess limits through several operationally relevant dimensions.


\begin{table}[htbp]\centering\small
\caption{Accuracy and hit diagnostics for explicitly named reports. Panel A reports paired working-reference tests; Panel B uses the fixed nonoverlap grid for hit tests. A better average score does not imply fewer breaches.}\label{tab:diagnostics}
\begin{tabular}{>{\raggedright\arraybackslash}p{4.3cm}>{\raggedright\arraybackslash}p{2.1cm}>{\raggedright\arraybackslash}p{2.1cm}>{\raggedright\arraybackslash}p{5.7cm}}\toprule
Panel A: paired contrast & DM p & HLN p & Named adjustment / p\\\midrule
FX original primary & 2.20e-07 & 2.31e-07 & Two-primary Bonferroni: 4.39e-07\\
FX matched history & 4.47e-09 & 4.83e-09 & Original matched Holm family: 2.86e-07\\
Treasury original primary & 0.00761 & 0.00939 & Original matched Holm family: 0.320\\
\bottomrule\end{tabular}
\par\smallskip\begin{minipage}{.98\linewidth}\footnotesize Original matched Holm uses 72 contrasts. The later regrouped matched family is separate; earlier-period writing-stage HAC comparisons have neither a new HLN field nor a new family adjustment.\end{minipage}

\par\medskip\begin{center}\small
\begin{tabular}{p{4.0cm}rrrr}\toprule
Panel B: report & Grid hits & Grid issues & Kupiec p & Christoffersen CC p\\\midrule
FX regression & 50 & 947 & 0.695 & 0.636\\
FX original rolling reference & 49 & 947 & 0.807 & 0.905\\
FX matched fixed reference & 49 & 947 & 0.807 & 0.646\\
Treasury regression & 7 & 40 & 0.00413 & 0.0136\\
Treasury historical reference & 4 & 40 & 0.199 & 0.262\\
\bottomrule\end{tabular}\end{center}
{\footnotesize Hit tests use the full fixed nonoverlap grid; the conditional-coverage likelihood conditions on its initial hit. Equality to the nominal tail rate is the stated diagnostic null, distinct from the population upper-quantile inequality. The grid does not create independence.}
\end{table}

\begin{figure}[htbp]\centering\includegraphics[width=.94\linewidth]{figures/v7-murphy-fx_hourly.pdf}
\caption{Murphy diagram for the original short-horizon FX primary. The full-breakpoint calculation finds crossings; the graph is a display of saved scores. Average pinball superiority does not imply dominance across elementary scores.}\label{fig:murphyfx}
\end{figure}

\begin{figure}[htbp]\centering\includegraphics[width=.94\linewidth]{figures/v7-murphy-nyfed_weekly.pdf}
\caption{Murphy diagram for the original Treasury-bill primary. The saved full-breakpoint comparison also crosses. A single average-score ordering leaves other consistent-score preferences unresolved.}\label{fig:murphybill}
\end{figure}

\FloatBarrier

\subsection{Controlled and within-record extensions}\label{sec:extensions}

The controlled mean-reverting level study with GARCH innovations rejects the proposed narrow-regime explanation \citep{bollerslev1986}. Each regression shows clear gains in only 2 of 12 target cells. Both also show clear gains in 6 of 60 outside cells, including homoskedastic and longer-horizon cases. The experiment uses 32 independent repetition streams with common innovations across conditions and 72 evaluation issues per repetition. Figure~\ref{fig:mechanism} reports relative mean losses with Monte Carlo standard errors. These specified processes therefore do not establish a universal history, volatility or horizon boundary.

\begin{figure}[htbp]\centering\includegraphics[width=.98\linewidth]{figures/v8-wp3-qr33-loss.pdf}
\caption{Relative mean loss of the original regression against matched historical headroom in the controlled level-persistence and volatility-clustering grid. Annotations give the relative change with its Monte Carlo standard error. Panels vary horizon and level persistence; cells vary history and variance persistence. Gains outside the prespecified region contradict a universal narrow-condition explanation.}\label{fig:mechanism}
\end{figure}

Same-dealer daily aggregation produces predominantly historical-rule score wins. Historical headroom wins 7 of the 8 specified comparisons. The exception is the longer-history four-predictor regression at the shorter daily horizon, with about a 3.0\% score reduction. It records 37 breaches versus 30 for historical headroom. Aggregation changes physical histories, horizons, update cadence, predictors, evaluation span and the sampled peak together. Source identity is held fixed, but this contrast does not identify a pure frequency effect.

An expanded fixed-parameter backcast preserves the reference-sensitive pattern but supplies no fresh independent validation. For original short-horizon FX regression, matched score falls about 12.2\%, while its rolling-reference advantage is about 5.3\%. Matched breaches are 416 for regression and 282 for historical headroom. The raw rolling-reference DM p-value is about 0.055 and remains descriptive. Earlier rolling-validation scoring already covered these target paths, and the smaller backcast is fully nested. Penalties were chosen later, so these overlapping windows cannot be interpreted as independent replications.

\FloatBarrier

\section{Implications for model validation}\label{sec:implications}

Peak-limit validation should begin with the information set and the comparison rule before examining a candidate model's favorable performance statistic. Admit a historical response only after its whole peak horizon has matured, and construct all transformations from the eligible training observations. Then compare the candidate with a historical rule using the same history and update schedule before returning to the original operational reference. This sequence separates the matched model comparison from the combined effects of a different history length and a different update schedule.


A practical validation record should also keep scores, breaches, and stability together rather than replacing one assessment with another after seeing the results. Report the mean quantile loss together with breach frequency, overrun amounts, and unused headroom under the same issue grid. Repeat the comparison across the available periods and related series, while displaying every failed or adverse configuration beside the successful settings. Table~\ref{tab:map} organizes observed comparison settings and the validation questions they raise, including the daily and longer-horizon exceptions.


For a limit manager, the key decision is whether extra conditioning improves the chosen monitoring objective reliably enough to justify the additional estimation. The original hourly FX result motivates further evaluation of its named regression, while the weekly and daily comparisons support demanding historical benchmarks. Calibration and combination comparisons should enter that decision through their own measured score and breach behavior, with their definitions kept explicit. The appropriate default is a matched validation process that lets comparative evidence, rather than model complexity, determine the preferred peak-limit report.


\begin{table}[htbp]\centering\small
\caption{Observed comparison settings and validation questions. Entries organize retrospective evidence rather than a causal frequency effect or deployment policy. A winning score leaves breach behavior and operating suitability to be assessed.}\label{tab:map}
\begin{tabular}{>{\raggedright\arraybackslash}p{3.3cm}>{\raggedright\arraybackslash}p{3.7cm}>{\raggedright\arraybackslash}p{4.0cm}>{\raggedright\arraybackslash}p{3.2cm}}\toprule
Frequency / history & Horizon & Evidence & Validation action\\\midrule
Hourly / longer history & Short & Repeated matched score gains; separate wrapper check & Compare named control and breaches\\
Hourly / shorter history & Short & Small score gain with more breaches & Require stable joint performance\\
Hourly / either history & Longer & Matched backcast gains coexist with reference-sensitive losses; coherence and preference sensitivity remain & Check the named reference, horizon order and loss preferences\\
Daily / same dealer & Horizons in nonempty days & Historical score wins predominate; one four-predictor gain has more breaches & Retain both controls and assess their own breaches\\
Weekly / short or long history & Short or longer & Frequent regression losses despite expanded penalties & Compare historical and parsimonious controls\\
Any source / endpoint bound & Specified horizon & Population premises and plug-in behavior differ & Check absolute displacement and dependence\\
\bottomrule\end{tabular}

\end{table}

\FloatBarrier

\section{Limitations}\label{sec:limitations}

The FX evidence comes from a single market maker's client-flow record, limiting conclusions about other dealers, currencies, and market-making arrangements. Its original evaluation history had already been viewed, and the earlier chronological comparison reuses familiar observations from the same source. These overlapping windows and the previously viewed later period describe stability within the record rather than independent dealer replication or prospective performance. The expanded prefix includes the earlier backcast and uses penalties selected after its issue dates. A stronger external assessment requires a compatible additional dealer source and a future evaluation protocol fixed before the new outcomes become available.


The weekly data are signed position-category aggregates, so offsetting exposures across dealers can reduce the aggregate while individual gross exposures remain large. Related categories share a reporting system and economic components, which limits an interpretation of their comparisons as independent replications across institutions. The downloaded economic-date series also differs from a reconstruction of the observations available at each historical publication time. The sampled target measures recorded positions; excursions between observations, settlement cash needs, and the costs of actual limit decisions require additional measurement.


The later analyses and working-reference tests describe the observed records after earlier results had informed the research process and additional comparisons. Their p-values therefore do not account for the full sequence of model, sample, diagnostic, and presentation choices made during that process. The finite calibration wrapper, empirical endpoint rule, and dependent financial paths also require distinctions from the population assumptions in Proposition~\ref{prop:levy}. The main evidence supports a bounded comparative forecasting result, with external validity, prospective calibration, and economic decision value remaining questions for future evaluation. The score uses a nineteen-to-one error-weight ratio for evaluation, which does not measure the actual costs of a limit decision.


\FloatBarrier

\section{Conclusion}\label{sec:conclusion}

Simple historical rules are strong competitors for predicting the largest sampled absolute position before a monitoring horizon ends. The original long-history FX regression gains against matched headroom survive in the broader backcast on previously examined observations, although score reductions coexist with more breaches. Longer-horizon backcast gains and the daily exception prevent a literal short-horizon or hourly-only interpretation. The controlled simulation also rejects the proposed narrow-regime restriction, while the same-dealer aggregation contrast leaves frequency mixed with measurement and clock changes. The defensible result is reference-sensitive comparative accuracy in these records, rather than a universal mechanism or validated operating-limit rule.


The weekly comparisons show that the successful FX specification does not transport reliably to related peak-position tasks at the lower frequency. Heavier regularization reduces the additional panel's median excess loss substantially, yet historical rules still win most comparisons and the conservative-selection exercise misses its success criteria. Lower average quantile loss, fewer breaches and stable rankings impose different requirements, and these data frequently rank them differently. For managers selecting peak limits, their joint assessment provides a firmer basis for action than a favorable score considered on its own.


Future work should evaluate these comparisons prospectively and obtain a compatible independent dealer record, while preserving the same response and information clock. A decision study could then measure the economic consequences of changing a limit, including monitoring actions, funding, hedging, and inventory adjustments. Such evidence would connect forecast quality to realized operating value under a specified policy rather than treating asymmetric error weights as observed costs. Peak-limit forecasting advances when its comparisons remain matched, its failures remain visible, and its practical claims stay tied to what the data measure.


\FloatBarrier

\section*{Data and code availability}
\paragraph{Data availability.} The FX input is the specified version of the Kaggle Net Position Time Series Dataset. The owning metadata labels its license ``Data files \copyright\ Original Authors.'' This repository supplies the acquisition route and input fingerprint, but excludes raw FX observations and private per-issue prediction vectors. Obtain the source from \url{https://www.kaggle.com/datasets/davidlindsay1979/net-position-time-series-dataset} subject to the original rights holders' terms.

The New York Fed aggregates are available from its Primary Dealer Statistics service under the Bank's Terms of Use. The repository preserves the unchanged accepted Treasury snapshot with attribution and identifies study calculations separately. Copying or distributing that content remains subject to its source notices, attribution, modification labelling, continuing conditions and endorsement restrictions: \url{https://www.newyorkfed.org/privacy/termsofuse.html}. The code license does not replace those source conditions.

\paragraph{Code availability.} Code, saved results and version-specific reproduction instructions are available at \url{https://github.com/ByronLi0207/financial-mathematics-paper-20261020}. This revision adds three design-frozen exploratory extensions: a controlled simulation, a same-dealer daily aggregation contrast and an expanded inspected backcast. The completed results include adverse, failed and unavailable cases; arithmetic readback establishes artifact consistency rather than independent scientific replication.

\section*{Author-provided information}
Author information: [TO BE COMPLETED BY AUTHOR].\par
ORCID: [TO BE COMPLETED BY AUTHOR].\par
Funding: [TO BE COMPLETED BY AUTHOR].\par
Competing interests: [TO BE COMPLETED BY AUTHOR].\par
AI statement: [TO BE COMPLETED BY AUTHOR].

\begin{thebibliography}{99}
\raggedright
\bibitem[Arjani et al.(2020)]{arjani2020} Arjani, N., Li, F., \& Sabetti, L. (2020). \emph{Monitoring intraday liquidity risks in a real time gross settlement system}. Payments Canada. \url{https://www.payments.ca/sites/default/files/discussion_paper_monitoring_intraday_liquidity_risks_in_a_real_time_gross_settlement_system.pdf}.

\bibitem[Barone-Adesi et al.(1999)]{barone1999} Barone-Adesi, G., Giannopoulos, K., \& Vosper, L. (1999). VaR without correlations for portfolios of derivative securities. \emph{Journal of Futures Markets}, \emph{19}(5), 583--602. \url{https://doi.org/10.1002/(SICI)1096-9934(199908)19:5<583::AID-FUT5>3.0.CO;2-S}.

\bibitem[Board of Governors of the Federal Reserve System(2013)]{fr2004_2013} Board of Governors of the Federal Reserve System (2013, March). \emph{Reporting guidelines for preparing the FR 2004 primary government securities dealers reports}. Board of Governors of the Federal Reserve System. \url{https://www.federalreserve.gov/apps/reportingforms/Download/DownloadAttachment?guid=f555fe11-e149-479b-8088-4ed4e9ae1206}.

\bibitem[Board of Governors of the Federal Reserve System(2024)]{fr2004_2024} Board of Governors of the Federal Reserve System (2024, July). \emph{Instructions for the preparation of government securities dealers reports: Reporting form FR 2004}. Board of Governors of the Federal Reserve System. \url{https://www.federalreserve.gov/apps/reportingforms/Download/DownloadAttachment?guid=950ab256-b485-4586-99b0-b27eb205fdfa}.

\bibitem[Bollerslev(1986)]{bollerslev1986} Bollerslev, T. (1986). Generalized autoregressive conditional heteroskedasticity. \emph{Journal of Econometrics}, \emph{31}(3), 307--327. \url{https://doi.org/10.1016/0304-4076(86)90063-1}.

\bibitem[Boudoukh et al.(2004)]{boudoukh2004} Boudoukh, J., Richardson, M., Stanton, R., \& Whitelaw, R. F. (2004). MaxVaR: Long-horizon value at risk in a mark-to-market environment. \emph{Journal of Investment Management}, \emph{2}(3), 1--6. \url{https://pages.stern.nyu.edu/~rwhitela/papers/maxvar%20joim04.pdf}.

\bibitem[Christoffersen(1998)]{christ1998} Christoffersen, P. F. (1998). Evaluating interval forecasts. \emph{International Economic Review}, \emph{39}(4), 841--862. \url{https://doi.org/10.2307/2527341}.

\bibitem[Diebold \& Mariano(1995)]{dm1995} Diebold, F. X., \& Mariano, R. S. (1995). Comparing predictive accuracy. \emph{Journal of Business \& Economic Statistics}, \emph{13}(3), 253--263. \url{https://doi.org/10.1080/07350015.1995.10524599}.

\bibitem[Duffie et al.(2023)]{duffie2023} Duffie, D., Fleming, M. J., Keane, F. M., Nelson, C., Shachar, O., \& Van Tassel, P. (2023). \emph{Dealer capacity and U.S. Treasury market functionality} (Staff Report No. 1070). Federal Reserve Bank of New York Staff Reports. \url{https://doi.org/10.59576/sr.1070}.

\bibitem[Ehm et al.(2016)]{ehm2016} Ehm, W., Gneiting, T., Jordan, A., \& Krüger, F. (2016). Of quantiles and expectiles: Consistent scoring functions, Choquet representations and forecast rankings. \emph{Journal of the Royal Statistical Society Series B: Statistical Methodology}, \emph{78}(3), 505--562. \url{https://doi.org/10.1111/rssb.12154}.

\bibitem[Engle \& Manganelli(2004)]{engle2004} Engle, R. F., \& Manganelli, S. (2004). CAViaR: Conditional autoregressive value at risk by regression quantiles. \emph{Journal of Business \& Economic Statistics}, \emph{22}(4), 367--381. \url{https://doi.org/10.1198/073500104000000370}.

\bibitem[Federal Reserve Bank of New York(n.d.)]{nyfeddata} Federal Reserve Bank of New York (n.d.). \emph{Primary dealer statistics} [Data set]. Federal Reserve Bank of New York. \url{https://www.newyorkfed.org/markets/counterparties/primary-dealers-statistics}. Retrieved October 23, 2026.

\bibitem[Giacomini \& White(2006)]{gw2006} Giacomini, R., \& White, H. (2006). Tests of conditional predictive ability. \emph{Econometrica}, \emph{74}(6), 1545--1578. \url{https://doi.org/10.1111/j.1468-0262.2006.00718.x}.

\bibitem[Gibbs \& Candès(2021)]{gibbs2021} Gibbs, I., \& Candès, E. J. (2021). Adaptive conformal inference under distribution shift. \emph{Advances in Neural Information Processing Systems}, \emph{34}, 1660--1672. \url{https://proceedings.neurips.cc/paper_files/paper/2021/hash/0d441de75945e5acbc865406fc9a2559-Abstract.html}.

\bibitem[Gneiting(2011)]{gneiting2011} Gneiting, T. (2011). Making and evaluating point forecasts. \emph{Journal of the American Statistical Association}, \emph{106}(494), 746--762. \url{https://doi.org/10.1198/jasa.2011.r10138}.

\bibitem[Harvey et al.(1997)]{hln1997} Harvey, D., Leybourne, S., \& Newbold, P. (1997). Testing the equality of prediction mean squared errors. \emph{International Journal of Forecasting}, \emph{13}(2), 281--291. \url{https://doi.org/10.1016/s0169-2070(96)00719-4}.

\bibitem[Holm(1979)]{holm1979} Holm, S. (1979). A simple sequentially rejective multiple test procedure. \emph{Scandinavian Journal of Statistics}, \emph{6}(2), 65--70. \url{https://www.jstor.org/stable/4615733}.

\bibitem[Hull \& White(1998)]{hull1998} Hull, J., \& White, A. (1998). Incorporating volatility updating into the historical simulation method for value-at-risk. \emph{The Journal of Risk}, \emph{1}(1), 5--19. \url{https://doi.org/10.21314/jor.1998.001}.

\bibitem[Koenker \& Bassett(1978)]{koenker1978} Koenker, R., \& Bassett, G., Jr. (1978). Regression quantiles. \emph{Econometrica}, \emph{46}(1), 33--50. \url{https://doi.org/10.2307/1913643}.

\bibitem[Kupiec(1995)]{kupiec1995} Kupiec, P. H. (1995). Techniques for verifying the accuracy of risk measurement models. \emph{The Journal of Derivatives}, \emph{3}(2), 73--84. \url{https://doi.org/10.3905/jod.1995.407942}.

\bibitem[Lindsay(2020)]{dataset} Lindsay, D. (2020). \emph{Net position time series dataset} (Version 2) [Data set]. Kaggle. \url{https://www.kaggle.com/datasets/davidlindsay1979/net-position-time-series-dataset}.

\bibitem[McNeil \& Frey(2000)]{mcneilfrey2000} McNeil, A. J., \& Frey, R. (2000). Estimation of tail-related risk measures for heteroscedastic financial time series: An extreme value approach. \emph{Journal of Empirical Finance}, \emph{7}(3--4), 271--300. \url{https://doi.org/10.1016/s0927-5398(00)00012-8}.

\bibitem[Meinshausen(2006)]{meinshausen2006} Meinshausen, N. (2006). Quantile regression forests. \emph{Journal of Machine Learning Research}, \emph{7}(35), 983--999. \url{https://www.jmlr.org/papers/v7/meinshausen06a.html}.

\bibitem[Newey \& West(1987)]{nw1987} Newey, W. K., \& West, K. D. (1987). A simple, positive semi-definite, heteroskedasticity and autocorrelation consistent covariance matrix. \emph{Econometrica}, \emph{55}(3), 703--708. \url{https://doi.org/10.2307/1913610}.

\bibitem[Petrunia et al.(2019)]{petrunia2019} Petrunia, R., Sabetti, L., \& Voia, M. (2019). \emph{Tail risk in the automated clearing settlement system (ACSS)} (Discussion Paper No. 7). Payments Canada. \url{https://www.payments.ca/sites/default/files/tail_risk_in_the_automated_clearing_settlement_system_acss.pdf}.

\bibitem[Sussman(2023)]{sussman2023} Sussman, E. (2023). A strengthened Orlicz--Pettis theorem via Itô--Nisio. \emph{Annals of Functional Analysis}, \emph{14}(1), Article 22. \url{https://doi.org/10.1007/s43034-022-00225-1}.

\bibitem[van den Hoogen(2017)]{hoogen2017} van den Hoogen, S. (2017). \emph{Intraday liquidity risk estimation using transaction data: An extreme value theory approach} [Master's thesis, Delft University of Technology]. \url{https://resolver.tudelft.nl/uuid:1c0c4545-bb07-4916-aa37-1d4bea3b7872}.

\bibitem[Wisniewski et al.(2020)]{wisniewski2020} Wisniewski, W., Lindsay, D., \& Lindsay, S. (2020). Application of conformal prediction interval estimations to market makers' net positions. In \emph{Proceedings of the Ninth Symposium on Conformal and Probabilistic Prediction and Applications} (Proceedings of Machine Learning Research, Vol. 128, pp. 285--301). \url{https://proceedings.mlr.press/v128/wisniewski20a.html}.

\bibitem[Wisniewski(2024)]{wisniewski2024} Wisniewski, W. J. (2024). \emph{Net position and retail trader behaviour: Implications for market makers' risk management strategies} [Doctoral thesis, Royal Holloway, University of London]. \url{https://pure.royalholloway.ac.uk/en/publications/net-position-and-retail-trader-behaviour-implications-for-market-/}.
\end{thebibliography}

\end{document}
